\abstract{We present a conservative Darcy-Brinkman-Stokes (DBS) formulation for
PHREEQC-coupled reactive transport with cell-based surface complexation in
image-based porous media. Implemented in the OpenFOAM-based GeoChemFoam
framework, the formulation transports dissolved species in the physical
mobile-water fraction, stores adsorbed species as conservative cell-scale
surface inventories, and couples each reactive cell to PHREEQC. Reactive
surface area is supplied independently as a cell field, allowing unresolved
internal area in segmented microporous phases to be represented while DBS flow
and transport are retained in resolved pore space.

The implementation is assessed through tests that separate conservative
transport--chemistry coupling from surface-area construction. A one-dimensional
analytical retardation benchmark and closed-system batch tests evaluate the
aqueous and adsorbed mass balance. Controlled two- and three-dimensional
geometries test the smooth-interface surface-area algorithms, while image-based
fracture-network and Estaillades limestone cases examine reactive transport in
heterogeneous pore structures.

The application set spans Pb$^{2+}$ retardation by hydrous ferric oxide
coatings in a two-dimensional fracture network, competitive Mn$^{2+}$/Mg$^{2+}$
exchange in a three-dimensional Kelvin-cell sorbent geometry, and
Sr$^{2+}$/Ca$^{2+}$ exchange in Estaillades limestone. The fracture-network and
limestone cases illustrate how the formulation produces spatially resolved
aqueous and adsorbed
inventories when reactive surface area is partly unresolved at the image scale.
The resulting framework supports conservative surface-complexation modelling
across resolved and unresolved pore structures relevant to contaminant
retention, selective ion capture, and radionuclide transport.}
